\begin{textsample}{NEG Dim 1 – 2022\_c4 – Score -4.00 – 2022\_c4/enade\_2022\_tec\_gest\_tur\_eadd874aa1.txt}  \label{ex:f1_neg_172}
CAPÍTULO 5 PERCEPÇÃO DA PROVA Nas análises feitas neste capítulo, constam as percepções dos concluintes da Área de Turismo sobre a prova aplicada no Enade/2022.
Essas percepções foram mensuradas por meio de nove questões que avaliaram desde o grau de dificuldade da prova até o tempo gasto para concluí -la.
As percepções sobre a prova foram relacionadas ao desempenho dos estudantes, à Grande região de funcionamento do curso, à Categoria administrativa e ao tipo de Organização acadêmica da IES.
O questionário de percepção da prova encontra -se ao final do Anexo VIII, no qual \textbf{está} a reprodução do exame.
Na apresentação dos dados relativos às nove questões sobre as percepções a respeito da prova, o desempenho dos estudantes foi classificado em quatro quartos.
Para tanto, esse desempenho foi ordenado de forma ascendente.
O percentil 25, P25, também conhecido como primeiro quartil, é a nota de desempenho que deixa um quarto ( 25 \% ) dos valores observados abaixo e três quartos acima.
O quarto inferior de desempenho é composto pelas notas abaixo do primeiro quartil.
Já o percentil 75, P75, também conhecido como terceiro quartil, é o valor para o qual há três quartos ( 75 \% ) dos dados abaixo e um quarto acima dele.
O quarto superior de desempenho é composto pelas notas iguais ou acima do terceiro quartil.
O percentil 50, P50, também conhecido como mediana, é o valor que divide as notas em dois conjuntos de igual tamanho.
O segundo quarto inclui os valores entre o primeiro quartil ( P25 ) e a mediana.
O terceiro quarto contém os valores entre a mediana ( P50 ) e o terceiro quartil ( P75 ).
Vale ressaltar que percentis, quartis e medianas são pontos que, não obrigatoriamente, pertencem ao conjunto original de dados, ao passo que os quartos são subconjuntos dos dados originais.
Na Figura 5.1, apresenta -se uma ilustração dos quatro quartos descritos.
A seguir, serão apresentados gráficos com resultados selecionados, relativos às nove questões avaliadas por grupos de estudantes.
Nas barras dos gráficos, apresenta -se o percentual de estudantes que assinalaram uma das opções ou a soma das porcentagens daqueles que assinalaram duas ( ou três ) delas.
Por exemplo, para as questões 1 e 2, nos gráficos, é apresentada a porcentagem total de participantes que assinalaram as opções ( D ) “ Difícil ” e ( E ) Muito “ Difícil ”.
Em \textbf{cada} barra, foram assinalados, também, os extremos do intervalo de confiança de 95 \% como linhas verticais unidas por uma linha horizontal na forma da letra H maiúscula.
O estimador de um parâmetro com um \textbf{certo} nível de confiança ( por exemplo, 95,0 \% ) deve conter o parâmetro no intervalo de confiança em 95 \% das vezes.
Na comparação entre os estimadores dos parâmetros de duas classes de uma dada categoria ( por exemplo : Norte e Nordeste nas Grandes regiões, de primeiro e último quarto dentro de desempenho, de IES públicas e privadas, ou de universidades e Faculdades ), associados aos seus respectivos intervalos de confiança, diz -se que não há diferença estatisticamente significativa entre os parâmetros das duas categorias quando há interseção entre os intervalos de confiança, e que há diferença, se os intervalos de confiança são disjuntos ( para maiores informações vide Glossário ).
Nas Tabelas do Anexo V, são apresentados os valores absolutos e a distribuição percentual9 das alternativas válidas das nove questões, segundo o mesmo recorte de desempenho dos estudantes e da grande região de funcionamento do curso, de categoria administrativa e do tipo de organização acadêmica da IES. 5.1 GRAU DE DIFICULDADE DA PROVA Nesta seção são mostradas as percepções dos estudantes quanto ao grau de dificuldade da prova Enade/2022, e que foram analisadas considerando sua vinculação ao Componente de Formação Geral e ao Componente específico. 9 Reitera -se que uma das convenções para tabelas numéricas ( Subseção 2.2.3 do Capítulo 2 ) refere -se à possibilidade de, por questão de arredondamento, a soma das partes não resultar em 100 \%. 5.1.1 COMPONENTE DE FORMAÇÃO GERAL O Gráfico 5.1 apresenta as respostas obtidas para a Questão 1 : “ Qual o grau de dificuldade desta prova na parte de Formação Geral? ”.
De acordo com os dados do Gráfico 5.1, a Região Sul apresentou o maior percentual de estudantes que consideraram a prova “ Difícil ” ou “ Muito “ Difícil ”, com 23,5 \%, seguida da região Nordeste, com 19,5 \%.
Na Região Sudeste, a média foi de 12,9 \%, o menor percentual.
Nas demais regiões, os percentuais foram 16,7 \% ( Norte ) e 16,5 \% ( Centro-Oeste ).
O Gráfico 5.2 traz mais resultados sobre a questão tratada no gráfico anterior, porém por quartos de desempenho.
Os dados do Gráfico 5.2 mostram que o percentual de estudantes que consideraram o Componente de Formação Geral da prova “ Difícil ” ou Muito “ Difícil ” foi decrescente em função dos quartos de desempenho : 24,2 \%, no 1º quarto e 10,3 \%, no 4º quarto, o grupo de melhor desempenho na prova.
Nos quartos de desempenho intermediários, o percentual de estudantes que consideraram a prova “ Difícil ” ou Muito “ Difícil ” foi de 17,2 \%, no 2º quarto, e de 13,6 \%, no 3º quarto.
O Gráfico 5.3 traz outros resultados também em relação à mesma questão retratada no gráfico anterior, porém por organização acadêmica das IES.
Como mostram os dados do Gráfico 5.3, o percentual de estudantes que consideraram o Componente de Formação Geral da prova “ “ Difícil ” ou “ Muito Difícil ” foi maior entre os estudantes de Faculdades, alcançando 37,5 \%, enquanto entre os estudantes de Centros Universitários e Universidades ”, as Organizações Acadêmicas com as menores incidências, o percentual alcançou 16,9 \% e 15,7 \%, respectivamente.
O Gráfico 5.4 apresenta os dados por categoria administrativa das IES.
Os dados do Gráfico 5.4 mostram que o percentual de estudantes que consideraram o Componente de Formação Geral da prova “ “ Difícil ” ou “ Muito Difícil ” foi maior entre os estudantes de cursos de IES privadas ( 17,1 \% ) do que IES públicas ( 16 \% ). 5.1.2 COMPONENTE ESPECÍFICO O Gráfico 5.5 apresenta as respostas obtidas para a Questão 2 : “ Qual o grau de dificuldade desta prova na parte de Componente específico? ”.
De acordo com os dados do Gráfico 5.5, a Região Sul apresentou o maior percentual de estudantes que consideraram o Componente específico da prova “ “ Difícil ” ou “ Muito “ Difícil ”, com 26,8 \%, seguida da região Centro-Oeste com 20,7 \%.
Na Região Sudeste, a média foi de 17,3 \% ( o \textbf{menor} percentual ), na Norte, 18,1 \%, e na Região Nordeste, 17,8 \%.
O Gráfico 5.6 traz mais resultados sobre a questão tratada no gráfico anterior, porém por quartos de desempenho.
Os dados do Gráfico 5.6 mostram que o percentual de estudantes que consideraram o Componente específico da prova “ “ Difícil ” ou “ Muito Difícil ” foi decrescente em função dos quartos de desempenho : 22,8 \%, no 1º quarto e 16,1 \%, no 4º quarto, o grupo de melhor desempenho na prova.
Nos quartos de desempenho intermediários, o percentual de estudantes que consideraram a prova “ Difícil ” ou “ Muito Difícil ” foi de 19,5 \%, no 2º quarto, e de 16,5 \%, no 3º quarto.
O Gráfico 5.7 traz outros resultados também em relação à mesma questão retratada no gráfico anterior, porém por organização acadêmica das IES.
Como mostram os dados do Gráfico 5.7, o percentual de estudantes que consideraram o Componente de Formação Geral da prova “ Difícil ” ou “ Muito Difícil ” foi maior entre os estudantes de Faculdades, alcançando 25 \%, enquanto entre os estudantes das Universidades, a Organização acadêmica com a menor incidência, o percentual alcançou 18 \%.
O Gráfico 5.8 apresenta os dados por categoria administrativa das IES.
Os dados do Gráfico 5.8 mostram que o percentual de estudantes que consideraram o Componente específico da prova “ Difícil ” ou “ Muito Difícil ” foi ligeiramente maior entre os estudantes de cursos de IES privadas ( 19,1 \% ).
Enquanto o percentual registrado em IES públicas foi de 18,6 \%. 5.2 EXTENSÃO DA PROVA EM RELAÇÃO AO TEMPO TOTAL O Gráfico 5.9 apresenta as respostas obtidas para a Questão 3 : “ Considerando a extensão da prova, em relação ao tempo total, você considera que a prova foi longa ou muito longa? ”.
Em relação aos dados do Gráfico 5.9, entre as Grandes regiões, os que avaliaram a prova como longa ou muito longa, em relação ao tempo total destinado à sua resolução, atingiram uma porcentagem que variou de 32,8 \%, na região Norte, até 45 \%, na região Centro-Oeste.
As demais regiões apresentaram 38,3 \% ( Sudeste ), 38,7 \% ( Nordeste ) e 44,3 \% ( Sul ).
O Gráfico 5.10 traz mais resultados sobre a questão tratada no gráfico anterior, porém por quartos de desempenho.
No Gráfico 5.10, observa -se que o percentual de estudantes que consideraram a prova longa ou muito longa foi menor no 4º quarto de desempenho ( 36,5 \% ), e maior, no 1º quarto inferior ( 45,9 \% ).
Os percentuais intermediários ficaram entre os quartos, 2º e 3º, 37,2 \% e 38,1 \%, respectivamente.
O Gráfico 5.11 traz outros resultados também em relação à mesma questão retratada no gráfico anterior, porém por organização acadêmica das IES.
De acordo com os dados do Gráfico 5.11, considerando -se a extensão da prova, o percentual de estudantes que avaliaram a prova como “ longa ” ou “ muito longa ”, em relação ao tempo total destinado à sua resolução, foi maior entre os estudantes de Faculdades, alcançando 56,2 \%, enquanto entre os estudantes de Universidades, a organização acadêmica com a menor incidência, esse índice atingiu 38,8 \%.
Entre os participantes de CEFET/IF, o índice atingiu 42,3 \% e, entre os de centros universitários, 41,5 \%.
O Gráfico 5.12 apresenta os dados por categoria administrativa das IES.
Os dados do Gráfico 5.12 mostram que o percentual de estudantes que consideraram a prova “ longa ” ou “ muito longa ”, em relação ao tempo total destinado à sua resolução, foi maior entre os estudantes de cursos de IES privadas ( 43,7 \% ) do que entre os estudantes de IES públicas ( 38 \% ). 5.3 COMPREENSÃO DOS ENUNCIADOS DAS QUESTÕES Também organizados em relação ao Componente de Formação Geral e ao Componente específico, os dados apresentados nesta seção dizem respeito à forma como os estudantes compreenderam os enunciados das questões da prova, considerando as formulações expressas nesses mesmos enunciados. 5.3.1 COMPONENTE DE FORMAÇÃO GERAL O Gráfico 5.13 apresenta as respostas obtidas para a Questão 4 : “ Os enunciados das questões da prova na parte de Formação Geral \textbf{estavam} claros e objetivos? ” Os resultados apresentados no Gráfico 5.13 mostram que, na análise regional, a percentagem de estudantes que avaliaram que “ todos ” ou “ a maioria ” dos enunciados das questões do Componente de Formação Geral \textbf{estavam} “ claros e objetivos ” variou entre 78,7 \%, na região Norte, e 84,3 \%, na região Sudeste.
O Gráfico 5.14 traz mais resultados sobre a questão tratada no gráfico anterior, porém por quartos de desempenho.
De acordo com os dados do Gráfico 5.14, segundo o desempenho, observa -se que o percentual dos que emitiram essa opinião cresce, conforme o desempenho aumenta.
No quarto superior, a clareza e a objetividade de “ todos ” ou da “ maioria ” dos enunciados das questões foram percebidas por 91,3 \% dos respondentes, enquanto no quarto inferior esse percentual foi de 69,5 \% O Gráfico 5.15 traz outros resultados também em relação à mesma questão retratada no gráfico anterior, porém por organização acadêmica das IES.
Os dados do Gráfico 5.15 mostram que o percentual dos presentes que consideraram que “ todos ” ou “ a maioria ” dos enunciados das questões do Componente de Formação Geral \textbf{estavam} “ claros e objetivos ” foi maior entre os estudantes de Universidades, chegando a 82,3 \%, enquanto entre os de Faculdades, a organização acadêmica com a menor incidência, esse índice atingiu 62,5 \%.
Entre os estudantes de Centros Universitários, o índice atingiu 73,8 \%, e, entre os de CEFET/IF, 81,4 \%.
O Gráfico 5.16 apresenta os dados por categoria administrativa das IES.
Os dados do Gráfico 5.16 mostram que o percentual de estudantes que consideraram que “ todos ” ou “ a maioria ” dos enunciados das questões do Componente de Formação Geral \textbf{estavam} “ claros e objetivos ” foi maior entre os estudantes de cursos de IES públicas ( 82,9 \% ), quando comparado aos estudantes de IES privadas ( 77,7 \% ). 5.3.2 COMPONENTE ESPECÍFICO O Gráfico 5.17 apresenta as respostas obtidas para a Questão 5 : “ Os enunciados das questões da prova na parte de Componente específico \textbf{estavam} claros e objetivos? ”.
A maioria dos estudantes de todas as grandes regiões brasileiras considerou “ claros e objetivos ” todos ou a maioria dos enunciados das questões do Componente específico da prova, com percentual sendo maior na região Sul ( 83,9 \% ), ao contrário da região Centro-Oeste, com \textbf{menor} percentual, 78,9 \%.
O Gráfico 5.18 traz mais resultados sobre a questão tratada no gráfico anterior, porém por quartos de desempenho.
Os dados do Gráfico 5.18 mostram que o percentual de estudantes que consideraram os enunciados das questões “ claros e objetivos ” apresenta uma tendência crescente em relação ao desempenho : mais elevada no quarto superior ( 86,2 \% ), se comparada ao quarto inferior de desempenho ( 74,4 \% ).
O Gráfico 5.19 traz outros resultados também em relação à mesma questão retratada no gráfico anterior, porém por organização acadêmica das IES.
Os dados do Gráfico 5.19 mostram que o percentual dos presentes que consideraram que todos ou a maioria dos enunciados das questões do Componente específico \textbf{estavam} “ claros e objetivos ” foi maior entre os estudantes de Universidades, atingindo 81,5 \%.
Entre os estudantes de Centros Universitários e CEFET/IF, ambas Organizações Acadêmicas com a menor incidência, esse percentual atingiu 80 \% e 80,8 \%, respectivamente, sendo intermediária o percentual alcançado entre os estudantes de Faculdades ( 81,2 \% ).
O Gráfico 5.20 apresenta os dados por categoria administrativa das IES.
Os dados do Gráfico 5.20 mostram que o percentual de estudantes que consideraram que “ todos ” ou “ a maioria ” dos enunciados das questões do Componente específico \textbf{estavam} “ claros e objetivos ” foi maior entre os estudantes de cursos de IES públicas ( 82 \% ) do que entre os de IES privadas ( 79,2 \% ). 5.4 SUFICIÊNCIA DAS INFORMAÇÕES/INSTRUÇÕES FORNECIDAS O Gráfico 5.21 apresenta as respostas obtidas para a Questão 6 : “ As informações/instruções fornecidas para a resolução das questões foram suficientes para resolvê -las? ” De acordo com os dados do Gráfico 5.21 e quanto à distribuição de respondentes pelas grandes regiões, observa -se que o percentual de estudantes que consideraram as informações/instruções fornecidas “ até excessivas ” ou “ suficientes ” em “ todas ” ou “ na maioria ” das questões foi maior na região Sul, com 92,6 \%, e o menor percentual foi registrado na região Norte, chegando a 84,3 \%.
O Gráfico 5.22 traz resultados sobre a questão tratada no gráfico anterior, porém por quartos de desempenho.
De acordo com os dados do Gráfico 5.22, levando -se em conta o desempenho dos participantes, observa -se uma tendência crescente dos percentuais de participantes que avaliaram as informações/instruções como “ até excessivas ” ou “ suficientes ” em “ todas ” ou “ na maioria ” das questões.
O percentual foi mais elevado no quarto superior ( 95 \% ).
Já no primeiro quarto inferior, a suficiência das informações/instruções declarada como até excessiva, em todas ou na maioria das questões, foi percebida por 83,6 \% dos respondentes.
O Gráfico 5.23 traz outros resultados também em relação à mesma questão retratada no gráfico anterior, porém por organização acadêmica das IES.
De acordo com os dados do Gráfico 5.23, considerando -se as informações/instruções fornecidas para a resolução das questões, o percentual que considerou que essas eram “ até excessivas ” ou “ suficientes ” em “ todas ” ou “ na maioria ” das questões foi maior entre os estudantes de Centros Universitários, totalizando 93,8 \%.
As Faculdades foram a Organização acadêmica em que houve o menor percentual : 68,8 \%.
As Universidades e CEFET/IF corresponderam, respectivamente, a 89,3 \% e 86,5 \%.
O Gráfico 5.24 apresenta os dados por categoria administrativa das IES.
De acordo com os dados do Gráfico 5.24, o percentual de estudantes que consideraram que as informações/instruções fornecidas para a resolução das questões eram “ até excessivas ” ou “ suficientes ” em “ todas ” ou “ na maioria ” das questões foi maior entre os estudantes de cursos de IES públicas ( 89,3 \% ), quando comparado a IES privadas ( 88,1 \% ). 5.5 DIFICULDADE ENCONTRADA AO RESPONDER À PROVA O Gráfico 5.25 apresenta as respostas obtidas para a Questão 7 : “ Você se deparou com alguma dificuldade ao responder à prova?
Qual? ” Nos Gráficos 5.25 a 5.28, apresentados de forma sequenciada mais adiante, são expostos os percentuais de estudantes que apontaram o “ Desconhecimento do conteúdo ” como “ dificuldade ao responder à prova ”.
O Gráfico 5.25 apresenta as respostas obtidas sobre o “ Desconhecimento do conteúdo ” como a principal “ dificuldade ao responder à prova ”.
De acordo com os dados o Gráfico 5.25, o percentual de inscritos e presentes que apontaram o “ Desconhecimento do conteúdo ” como “ dificuldade ao responder à prova ” foi maior na região Sul com 12,2 \% e menor nas regiões Norte e Sudeste com 10,3 \% em \textbf{cada} uma.
O Gráfico 5.26 traz resultados sobre o ponto tratado no gráfico anterior, porém por quartos de desempenho.
De acordo com os dados do Gráfico 5.26, em relação aos quartos de desempenho, o “ Desconhecimento do conteúdo ” foi a opção escolhida por 12,2 \% dos estudantes do 3º quarto superior e por 8,6 \% no 2º quarto inferior.
O primeiro e último quarto obtiveram, respectivamente, 11,5 \% e 11 \%.
O Gráfico 5.27 traz outros resultados também em relação à mesma questão retratada no gráfico anterior, porém por organização acadêmica das IES.
De acordo com os dados do Gráfico 5.27, na análise por Organização acadêmica, o percentual de inscritos e presentes que apontaram o “ Desconhecimento do conteúdo ” como “ dificuldade ao responder à prova ” foi de 31,2 \% para estudantes de Faculdades ( o maior percentual ).
Já os estudantes de Centros Universitários apresentaram o \textbf{menor} percentual, 1,6 \%.
As Universidades e CEFET/IF apresentaram valores intermediários bem próximos, 10,9 \% e 11,6 \%, respectivamente.
O Gráfico 5.28 apresenta os dados por categoria administrativa das IES.
De acordo com os dados do Gráfico 5.28, em relação às categorias administrativas, o “ Desconhecimento do conteúdo ” foi a opção escolhida por 10,8 \% dos estudantes, nas IES públicas, e nas IES privadas, 10,9 \%. 5.6 CONTEÚDOS DAS QUESTÕES OBJETIVAS DA PROVA O Gráfico 5.29 apresenta as respostas obtidas para a Questão 8 : “ Considerando apenas as questões objetivas da prova, você percebeu que não estudou ainda a maioria desses conteúdos? ”.
Na análise por grande região, o percentual de respondentes que escolheram a opção “ não estudou ainda a maioria desses conteúdos ” foi pequena.
Observa -se que, nas regiões Sul ( 3,4 \% ) e Norte ( 3,1 \% ), os percentuais foram ligeiramente maiores.
Na região Sudeste, a situação foi inversa, com percentual mais baixo, correspondendo a 1,9 \%.
As regiões Nordeste e Centro-Oeste obtiveram o mesmo percentual intermediário, 2,8 \%.
O Gráfico 5.30 traz resultados sobre a questão tratada no gráfico anterior, porém por quartos de desempenho.
De acordo com os dados do Gráfico 5.30, considerando -se separadamente as opiniões de estudantes dos quatro quartos de desempenho, observa -se que, há uma tendência decrescente, ou seja, no quarto inferior, 5,3 \% ofereceram como resposta que “ não estudaram ainda a maioria desses conteúdos ”, sendo 0,9 \% os do quarto superior com a mesma resposta.
O Gráfico 5.31 traz outros resultados também em relação à mesma questão retratada no gráfico anterior, porém por organização acadêmica das IES.
De acordo com os dados do Gráfico 5.31, na análise por Organização acadêmica, o percentual de respondentes que escolheram a opção “ não estudou ainda a maioria desses conteúdos ” foi pequena.
Observa -se que os percentuais nos CEFET/IF ( 2,6 \% ) e nas Universidades ( 2,3 \% ) foram menores.
Nas Faculdades, a situação foi inversa : o percentual foi mais alto, 12,5 \%.
Os Centros Universitários ficaram no nível intermediário, 3,1 \%.
O Gráfico 5.32 apresenta os dados por categoria administrativa das IES.
De acordo com os dados do Gráfico 5.32, em relação às categorias administrativas, observa -se que, nas IES públicas, 2,2 \% ofereceram como resposta que “ não estudou ainda a maioria desses conteúdos ”, sendo 3,2 \% nas IES privadas com a mesma resposta. 5.7 TEMPO GASTO PARA CONCLUIR A PROVA O Gráfico 5.33 apresenta as respostas obtidas para a Questão 9 : “ Qual foi o tempo gasto por você para concluir a prova? ” Considerando -se as cinco grandes regiões brasileiras, os percentuais dos que utilizaram entre duas e quatro horas para finalizar a prova foram maiores, nas regiões Nordeste e Centro-Oeste, respectivamente, 85,1 \% e 84,8 \%.
Nas regiões Sudeste e Norte, os percentuais de estudantes que utilizaram entre duas e quatro horas para concluir a prova ficaram, 75,1 \% e 79,2 \%, respectivamente, apresentando os menores percentuais.
A região Sul ficou com o percentual de 82,6 \%.
O Gráfico 5.34 traz resultados sobre a questão tratada no gráfico anterior, porém por quartos de desempenho.
De acordo com os dados do Gráfico 5.34, uma vez analisadas as alternativas escolhidas pelos estudantes que se situam nos diferentes quartos de desempenho, observam -se uma tendência crescente : um percentual maior de participantes no quarto superior ( 85,7 \% ) declarou ter gasto entre duas e quatro horas para concluir a prova quando comparada com o percentual de participantes no primeiro quarto inferior ( 70,9 \% ).
O Gráfico 5.35 traz outros resultados também em relação à mesma questão retratada no gráfico anterior, porém por organização acadêmica das IES.
De acordo com os dados do Gráfico 5.35, na análise por Organização acadêmica, os percentuais dos que utilizaram entre duas e quatro horas para finalizar a prova foram de 90,8 \%, nos Centros Universitários ; 87,7 \%, nos CEFET/IF e 93,8 \%, nas Faculdades ( maior índice ).
O menor percentual foi entre os estudantes de Universidades, com 78,4 \% dos alunos necessitando de duas a quatro horas para concluir a prova.
O Gráfico 5.36 apresenta os dados por categoria administrativa das IES.
De acordo com os dados do Gráfico 5.36, em relação às categorias administrativas, observa -se que, nas IES públicas, 81,8 \% declararam ter gasto entre duas e quatro horas para concluir a prova, sendo 73,1 \% nas IES privadas, com a mesma resposta.
Os resultados do Enade 2022 apresentados neste capítulo contribuem para o entendimento de que a consideração da percepção dos estudantes concluintes sobre a prova do Enade é importante tanto para a compreensão do seu desempenho como para a reflexão sistemática e constante a respeito de determinadas realidades relacionadas a esse mesmo desempenho, no país.
Por outro lado, como todos os elementos do trabalho pedagógico institucional mostram -se diretamente relacionados com o desempenho dos estudantes, também ganha relevância conhecer e analisar os resultados da percepção dos coordenadores de curso sobre a mesma prova.
As informações expressas por estudantes e coordenadores sobre questões diversas associadas aos dois componentes da prova – Formação Geral e Conhecimento Específico – podem agregar valor em ações institucionais voltadas para o aperfeiçoamento do instrumento avaliativo, considerando os vários elementos envolvidos em sua elaboração.
Neste contexto, é relevante que os resultados das percepções de estudantes e coordenadores sobre a prova do Enade ganhem peso como um elemento a mais nas análises da oferta dos cursos nas modalidades presencial e a distância, em nível nacional.
Espera -se, assim, que, em uma visão formativa de avaliação, o conhecimento e a análise desses resultados possam vir a contribuir para a construção de alternativas institucionais que venham a induzir a qualidade da educação superior.

% matched lemmas: cada, certo, estar, menor
\end{textsample}
